\documentclass[preview]{standalone}

\usepackage{amsmath}
\usepackage{amssymb}
\usepackage{stellar}
\usepackage{definitions}
\usepackage{bettelini}

\begin{document}

\id{linearalgebra-euclidean-affine-spaces}
\genpage

\section{Euclidean affine spaces}

\begin{snippetdefinition}{euclidean-affine-space-definition}{Euclidean affine space}
    A \emph{Euclidean affine space} is a real \affinespace
    \(\mathcal{A}\) whose direction \vectorspace \(V\) is endowed with an
    \innerproduct
    \[
        \langle\cdot,\cdot\rangle.
    \]
\end{snippetdefinition}

\begin{snippet}{euclidean-affine-space-intuition}
    The \innerproduct is not placed on the points of \(\mathcal{A}\), but on
    its direction \vectorspace. This is enough to measure lengths of vectors
    \(\vec{pq}\), and therefore distances between points \(p,q\in
    \mathcal{A}\), without choosing an origin.
\end{snippet}

\begin{snippetdefinition}{affine-euclidean-distance-definition}{Affine Euclidean distance}
    Let \(\mathcal{A}\) be a \euclideanaffinespace with direction
    \innerproductspace \(V\). The \emph{Euclidean distance} between
    \(p,q\in\mathcal{A}\) is
    \[
        d(p,q)=\|\vec{pq}\|.
    \]
\end{snippetdefinition}

\begin{snippetproposition}{affine-euclidean-distance-is-metric}{The affine Euclidean distance is a metric}
    Let \(\mathcal{A}\) be a \euclideanaffinespace. The function
    \[
        d\colon\mathcal{A}\cartesianprod\mathcal{A}\fromto\realnumbers,
        \qquad
        d(p,q)=\|\vec{pq}\|
    \]
    is a distance function on \(\mathcal{A}\).
\end{snippetproposition}

\begin{snippetproof}{affine-euclidean-distance-is-metric-proof}{affine-euclidean-distance-is-metric}{The affine Euclidean distance is a metric}
    For every \(p,q\in\mathcal{A}\),
    \[
        d(p,q)=\|\vec{pq}\|\geq 0.
    \]
    Also,
    \[
        d(p,q)=0
        \Longleftrightarrow
        \|\vec{pq}\|=0
        \Longleftrightarrow
        \vec{pq}=0
        \Longleftrightarrow
        p=q.
    \]
    Symmetry follows from
    \[
        d(q,p)=\|\vec{qp}\|=\|-\vec{pq}\|=\|\vec{pq}\|=d(p,q).
    \]
    Finally, by the Chasles relation and the triangle inequality for the
    norm,
    \[
        d(p,r)
        =
        \|\vec{pr}\|
        =
        \|\vec{pq}+\vec{qr}\|
        \leq
        \|\vec{pq}\|+\|\vec{qr}\|
        =
        d(p,q)+d(q,r).
    \]
\end{snippetproof}

\section{Affine isometries}

\begin{snippetdefinition}{affine-isometry-definition}{Affine isometry}
    Let \(\mathcal{A}\) and \(\mathcal{B}\) be
    \euclideanaffinespace[Euclidean affine spaces]. An \affinemap
    \[
        \varphi\colon\mathcal{A}\fromto\mathcal{B}
    \]
    is an \emph{affine isometry} if it preserves distances:
    \[
        d(\varphi(p),\varphi(q))=d(p,q)
    \]
    for every \(p,q\in\mathcal{A}\).
\end{snippetdefinition}

\begin{snippettheorem}{affine-isometry-linear-part-criterion}{Criterion through the linear part}
    Let \(\mathcal{A}\) and \(\mathcal{B}\) be
    \euclideanaffinespace[Euclidean affine spaces], with direction
    \innerproductspace[inner product spaces] \(V\) and \(W\). Let
    \[
        \varphi\colon\mathcal{A}\fromto\mathcal{B}
    \]
    be an \affinemap. Then \(\varphi\) is an \affineisometry if and only if
    its \associatedlinearmap
    \[
        \vec{\varphi}\colon V\fromto W
    \]
    is a \linearisometry, equivalently
    \[
        \|\vec{\varphi}(v)\|=\|v\|
    \]
    for every \(v\in V\).
\end{snippettheorem}

\begin{snippetproof}{affine-isometry-linear-part-criterion-proof}{affine-isometry-linear-part-criterion}{Criterion through the linear part}
    For every \(p,q\in\mathcal{A}\),
    \[
        \vec{\varphi(p)\varphi(q)}
        =
        \vec{\varphi}(\vec{pq}).
    \]
    Hence
    \[
        d(\varphi(p),\varphi(q))
        =
        \|\vec{\varphi(p)\varphi(q)}\|
        =
        \|\vec{\varphi}(\vec{pq})\|.
    \]
    Since
    \[
        d(p,q)=\|\vec{pq}\|,
    \]
    the map \(\varphi\) preserves distances if and only if
    \[
        \|\vec{\varphi}(\vec{pq})\|=\|\vec{pq}\|
    \]
    for every \(p,q\in\mathcal{A}\). Every \(v\in V\) can be written as
    \(v=\vec{pq}\) for suitable \(p,q\in\mathcal{A}\), so this condition is
    equivalent to
    \[
        \|\vec{\varphi}(v)\|=\|v\|
    \]
    for every \(v\in V\). By the characterization of
    \linearisometry[linear isometries], this is equivalent to
    \(\vec{\varphi}\) being a \linearisometry.
\end{snippetproof}

\begin{snippetcorollary}{canonical-affine-isometries-coordinate-criterion}{Affine isometries in canonical coordinates}
    In the canonical \euclideanaffinespace
    \[
        \mathbb{A}^n(\realnumbers),
    \]
    an \affinemap
    \[
        \varphi(x)=Ax+b
    \]
    is an \affineisometry if and only if \(A\) is an
    \orthogonalmatrix, equivalently
    \[
        A^\transpose A=\identmatrix{n}.
    \]
\end{snippetcorollary}

\begin{snippetproof}{canonical-affine-isometries-coordinate-criterion-proof}{canonical-affine-isometries-coordinate-criterion}{Affine isometries in canonical coordinates}
    The \associatedlinearmap of \(\varphi(x)=Ax+b\) is
    \[
        x\mapsto Ax.
    \]
    By the linear part criterion, \(\varphi\) is an \affineisometry if and
    only if this linear map is a \linearisometry of \(\realnumbers^n\). By
    the matrix criterion for \linearisometry[linear isometries], this is
    equivalent to
    \[
        A^\transpose A=\identmatrix{n}.
    \]
\end{snippetproof}

\begin{snippetcorollary}{translations-are-affine-isometries}{Translations are affine isometries}
    Let \(\mathcal{A}\) be a \euclideanaffinespace. Every
    \translationmap
    \[
        t_u\colon\mathcal{A}\fromto\mathcal{A}
    \]
    is an \affineisometry.
\end{snippetcorollary}

\begin{snippetproof}{translations-are-affine-isometries-proof}{translations-are-affine-isometries}{Translations are affine isometries}
    The \associatedlinearmap of a \translationmap is
    \(\identityfunc_V\). Since \(\identityfunc_V\) is a \linearisometry, the
    linear part criterion implies that the translation is an
    \affineisometry.
\end{snippetproof}

\end{document}
